\documentclass{report}
\usepackage{listings}
\usepackage{graphicx}

\begin{document}
\title{Report 3 - CUDA}
\maketitle

\section{Introduction}
The plan is to convert the picture to grayscale by replacing the R,G,B values of each pixel with their average.
I will try with different picture sizes to compare the gap.

\begin{itemize}
    \item Load an image from file (matplotlib’s imread)
    \item Flatten image into 1D array of RGB (reshape(pixelCount, 3))
    \item Implement grayscale using CPU (for range)
    \item Implement grayscale using GPU
    \item Save/show the image after each grayscale to validate the result visually, dtype=np.uint8
    \item Use time.time() to measure speedup
\end{itemize}

\section{Work}

\subsection{Importing data}
Importing data, checking the size and showing it with pyplot to ensure the import worked correctly.
\begin{lstlisting}[language=Python]
from numba import cuda
from matplotlib import pyplot

img = pyplot.imread("image.jpg")

height = len(img)
width = len(img[0])
pixelCount = height*width
print(height,"*",width," = ",pixelCount," pixels")

pyplot.imshow(img)
pyplot.show()
\end{lstlisting}
\begin{enumerate}
    \item Output: 108 * 120 = 12960 pixels         (13k)
    \item Output: 960 * 1280 = 1228800 pixels      (1.2M)
    \item Output: 3456 * 5184 = 17915904 pixels    (18M)
\end{enumerate}
\begin{figure}[htbp]
    \centering
    \includegraphics[width=0.4\linewidth]{frog.png}
    \includegraphics[width=0.4\linewidth]{frog3.png}
    \includegraphics[width=0.4\linewidth]{frog2.png}
\end{figure}

\subsection{Grayscale with CPU}
Importing the remaining missing dependencies. Implementing grayscale following this formula:
\begin{figure}
    \centering
    \includegraphics[width=0.5\linewidth]{grayscale_formula.png}
\end{figure}
I was having issues with the types (Python moment) and a uint8 overflow on the intermediary sum (in the average calculation) so I had to make (R,G,B) uint16 before converting the avg back to uint8 for my 8-bit RGB.
\begin{lstlisting}[language=Python]
import numpy as np
import time

#initializing the result array /w uint8 dtype (default=float)
gray_array = np.zeros((pixelCount, 3), np.uint8)

def grayscale_CPU(r_img):
  for i in range(pixelCount):
    #making the types real clear as an overflow was causing noise on the final image
    (r,g,b)=np.int16((r_img[i,0], r_img[i,1], r_img[i,2]))
    gray = np.uint8((r+g+b)/3)
    gray_array[i,0] = gray_array[i,1] = gray_array[i,2] = gray
  return gray_array

r_img = img.reshape(pixelCount, 3)

start=time.time()
gray_array = grayscale_CPU(r_img)
stop=time.time()
print("duration",stop-start)

gray_img = gray_array.reshape(height, width, 3)

pyplot.imshow(gray_img)
pyplot.show()
\end{lstlisting}
\begin{enumerate}
    \item Output: duration 0.04524660110473633 
    \item Output: duration 9.251688241958618
    \item Output: duration 70.6480085849762
\end{enumerate}

\begin{figure}
    \centering
    \includegraphics[width=0.4\linewidth]{frog_CPU.png}
    \includegraphics[width=0.4\linewidth]{frog3_CPU.png}
    \includegraphics[width=0.4\linewidth]{frog2_CPU.png}
\end{figure}

\subsection{Grayscale with GPU}
Based on Prof. implementation, added the round up of gridSize.
\begin{lstlisting}[language=Python]
@cuda.jit
def grayscale(src, dst):
  # where are we in the input?
  tidx = cuda.threadIdx.x + cuda.blockIdx.x * cuda.blockDim.x
  g = np.uint8((src[tidx, 0] + src[tidx, 1] + src[tidx, 2]) / 3)
  dst[tidx, 0] = dst[tidx, 1] = dst[tidx, 2] = g

devSrc = cuda.to_device(r_img)
devDst = cuda.device_array_like(r_img)
blockSize = 32

#rounding up
gridSize = (pixelCount + blockSize - 1) // blockSize

start=time.time()
#kernel call
grayscale[gridSize, blockSize](devSrc, devDst)
stop=time.time()
print("duration",stop-start)

gray_img = devDst.copy_to_host().reshape(height, width, 3)

pyplot.imshow(gray_img)
pyplot.show()
\end{lstlisting}
\begin{enumerate}
    \item Output: duration 0.542658805847168
    \item Output: duration 0.3926730155944824
    \item Output: duration 0.39334869384765625
\end{enumerate}

\begin{figure}
    \centering
    \includegraphics[width=0.4\linewidth]{frog_GPU.png}
    \includegraphics[width=0.4\linewidth]{frog3_GPU.png}
    \includegraphics[width=0.4\linewidth]{frog2_GPU.png}
\end{figure}

\section{Conclusion}
CPU vs GPU computation
\begin{enumerate}
    \item 0.04 vs 0.54
    \item 9.25 vs 0.39
    \item 70.64 vs 0.39
\end{enumerate}

The bigger the load (more pixels) the higher the gap between CPU and GPU speed.
On a small picture (13k pixels) the CPU was \textbf{10x faster} than the GPU.
However, on a big picture (18M pixels), the CPU was roughly \textbf{200x slower} than the GPU.
\end{document}